\section{Preliminaries}
\label{sec:prelim}

\subsection{Causal Effect of Anti-Seizure Medication}

Here we formalize the problem of estimating causal effects from observational longitudinal
data with time-varying confounding and treatment--confounder feedback loops
\cite{parikh2023effects,hernan2020causal,young2020causal,young2011comparative,wanis2024grace,robins2009estimation}.
We consider a longitudinal study observed at discrete times $t \in \{0,1,2,\ldots\}$.
Throughout, we use bar notation $\bar X_t$ to denote the history of a variable $X$ up to
time $t$.

\paragraph{Observed data.}
The variables involved in the analysis are:
\begin{itemize}
\item $L_t$: The scalar burden of epileptiform activity (EA) measured from EEG at time
$t$, as defined in Section~\ref{sec:notation}. This is the quantity the treatment policy
titrates against. We write $\bar L_t = (L_0,L_1,\ldots,L_t)$. For concreteness, $L_t$
may be operationalized as the seizure burden per monitoring epoch.

\item $X$: Static baseline clinical characteristics (e.g., age, SOFA score).

\item $H_t = (X,\bar L_t,\bar A_{t-1})$: The observed history through time $t$. This is
the conditioning set used in the identification assumptions below. Note the distinction
maintained throughout: policies are functions of the scalar $L_t$, whereas identification
conditions on the full history $H_t$.

\item $A_t$: Treatment administered at time $t$. In our clinical setting, this is a
continuous, nonnegative infusion rate (e.g., propofol infusion). We write
$\bar A_t = (A_0,A_1,\ldots,A_t)$.

\item $C_{t+1}$: Indicator of censoring during the interval $(t,t+1]$, where
$C_{t+1}=1$ if censoring occurs in $(t,t+1]$ and $C_{t+1}=0$ otherwise. Thus,
$\bar C_t=(C_1,\ldots,C_t)$ denotes the censoring history through interval $(t-1,t]$,
and the event ``$\bar C_t=0$'' means uncensored through time $t$.

\item $Y_{t+1}$: Outcome indicator during the interval $(t,t+1]$, where $Y_{t+1}=1$
if the event occurs in $(t,t+1]$ and $Y_{t+1}=0$ otherwise. We write
$\bar Y_t=(Y_1,\ldots,Y_t)$, so ``$\bar Y_t=0$'' means event-free through time $t$.

\item $Y^{\bar a,\bar c=0}_{t+1}$: The counterfactual outcome during $(t,t+1]$ under
treatment history $\bar a$ and hypothetical elimination of censoring (i.e., $\bar c=0$).
\end{itemize}

\paragraph{Time ordering.}
At each time $t$, treatment $A_t$ is assigned based on past observed history, and then
over the subsequent interval $(t,t+1]$ the outcome, censoring, and next covariate value
are realized. Formally, the within-interval ordering is
\begin{equation}
\{\bar L_t,\bar A_{t-1},\bar C_t,\bar Y_t\}
\longrightarrow A_t
\longrightarrow \{Y_{t+1},C_{t+1},L_{t+1}\}.
\end{equation}
These relationships are illustrated in the causal influence diagram in
Figure~\ref{fig:DAG2}.

\begin{figure}[H]
\centering
\begin{tikzpicture}[
  node distance = 1cm and 1.6cm,
  >=Latex, baseline,
  every node/.style={inner sep=0pt, anchor=center},
  main/.style={line width=0.9pt, draw=black},
  uniform/.style={line width=0.9pt, draw=black}
]
  % ----- time-0 nodes -----
  \node (L0) {$L_{0}$};
  \node[right=of L0] (A0) {$A_{0}$};
  \node[right=of A0] (Y0) {$Y_{0}$};
  \node[right=of Y0] (C1) {$C_{1}$};
  % ----- time-1 nodes (shifted right) -----
  \node[below=2cm of L0, xshift=1cm] (L1) {$L_{1}$};
  \node[right=of L1] (A1) {$A_{1}$};
  \node[right=of A1] (Y1) {$Y_{1}$};
  \node[right=of Y1] (C2) {$C_{2}$};
  % ----- within-time arrows (order inside each interval) -----
  \draw[->,uniform] (L0) -- (A0);
  \draw[->,uniform] (A0) -- (Y0);
  \draw[->,uniform] (Y0) -- (C1);
  \draw[->,uniform] (L1) -- (A1);
  \draw[->,uniform] (A1) -- (Y1);
  \draw[->,uniform] (Y1) -- (C2);
  % ----- cross-time "same-variable" persistence -----
  \draw[->,uniform,bend left=20] (L0) to (L1);
  \draw[->,uniform,bend left=20] (A0) to (A1);
  \draw[->,uniform,bend left=20] (Y0) to (Y1);
  \draw[->,uniform,bend left=20] (C1) to (C2);
  % ----- cross-time: all past → all future (now uniform black) -----
  % from L0
  \draw[->,uniform,bend left=16] (L0) to (A1);
  \draw[->,uniform,bend left=22] (L0) to (Y1);
  \draw[->,uniform,bend right=8] (L0) to (C2);
  % from A0
  \draw[->,uniform,bend left=12] (A0) to (L1);
  \draw[->,uniform,bend left=18] (A0) to (Y1);
  \draw[->,uniform,bend left=24] (A0) to (C2);
  % from Y0
  \draw[->,uniform,bend left=10] (Y0) to (L1);
  \draw[->,uniform,bend left=14] (Y0) to (A1);
  \draw[->,uniform,bend left=20] (Y0) to (C2);
  % from C1
  \draw[->,uniform,bend left=12] (C1) to (L1);
  \draw[->,uniform,bend left=18] (C1) to (A1);
  \draw[->,uniform,bend left=24] (C1) to (Y1);
\end{tikzpicture}
\caption{%
Causal diagram for a two-period longitudinal data structure illustrating time-varying
treatment, disease burden, censoring, and outcome processes. At each discrete time $t$,
the patient’s observed covariates $L_t$ summarize measured prognostic information
available prior to treatment assignment. Conditional on the history
$(\bar L_t, \bar A_{t-1})$, treatment $A_t$ is assigned by the clinical decision
process. Over the subsequent interval $(t,t+1]$, the patient may experience the outcome
$Y_{t+1}$ (e.g., death), may become censored $C_{t+1}$, and the disease process evolves
to $L_{t+1}$. Arrows indicate causal dependence across and within time periods,
including feedback between past treatment and future disease severity, motivating use
of the parametric g-formula.}
\label{fig:DAG2}
\end{figure}

\paragraph{Treatment strategies and estimands.}
A treatment strategy $g=(g_0,\ldots,g_K)$ is a deterministic dynamic rule that assigns
\[
a_t = g_t(\bar L_t,\bar a_{t-1}), \qquad t=0,\ldots,K.
\]
Under an intervention $g$, the treatment path $\bar a^g=(a_0^g,\ldots,a_K^g)$ is fixed
given the realized covariate history, even though the corresponding observed treatment
variables $\bar A$ are stochastic in observational data. Under the hypothetical
elimination of censoring, we write the counterfactual outcome at horizon $K+1$ as
\[
Y^g_{K+1} \equiv Y^{\bar a^g,\bar c=0}_{K+1}.
\]
Our goal is to compare risks under different strategies $g$ and $g'$ via contrasts such
as $\Pr(Y^g_{K+1}=1)-\Pr(Y^{g'}_{K+1}=1)$.

In our application, the primary strategy of interest $g$ is a closed-loop control policy
that mimics ICU titration of anti-seizure medication to maintain EA burden below a
target $\theta$ (Figure~\ref{fig:control-problem}D). A simple proportional--integral (PI)
policy can be written as
\begin{equation}
g_t(\bar L_t,\bar A_{t-1})
=
\min\!\left\{A_{\max},\,
\max\!\left\{0,\, \kappa_p e_t + \kappa_i \sum_{s=0}^{t} e_s \right\}\right\},
\qquad e_t = L_t-\theta,
\end{equation}
where $\kappa_p$ and $\kappa_i$ are proportional and integral gains and $A_{\max}$ is a
maximum safe dose. This closed-loop structure induces time-varying confounding: sicker
patients (higher $L_t$) tend to receive more treatment, while treatment affects future
$L_{t+1}$.

\subsection{Identification Assumptions}

Before using observational data to estimate causal effects, we specify assumptions under
which the counterfactual estimands are identifiable.

In plain language, three conditions must hold. (i) \emph{Sequential ignorability}: once we
condition on everything we have observed up to time $t$, the clinician's choice of dose at
$t$ and the chance of the patient leaving the study over the next interval are statistically
independent of what the patient's future survival would be under any given intervention.
This is an assumption about the \emph{assignment mechanism}, not a claim that treatment has
no effect. (ii) \emph{Positivity / overlap}: for histories that actually occur in the data,
the observed treatment and censoring mechanisms must have nonzero probability of producing
(or approaching, for continuous treatment) the action the intervention \(g\) would prescribe,
while the patient remains alive and uncensored. (iii) \emph{Consistency}: if the observed
treatment happens to match the intervention's prescription, the observed outcome equals the
counterfactual outcome under that intervention. Below we state these formally.

For each time $t=0,1,\ldots,K$, among individuals alive and uncensored at time $t$
(i.e., $\bar Y_t=\bar C_t=0$), we assume the following. Throughout, conditioning statements
are written in terms of $(\bar L_t,\bar A_{t-1})$; baseline covariates $X$ are fixed over
follow-up and are implicitly conditioned on in every such statement, so that the full
conditioning set is the observed history $H_t=(X,\bar L_t,\bar A_{t-1})$.

\textbf{(1) Sequential ignorability (no unmeasured confounding and no unmeasured informative censoring).}
Treatment and censoring behave ``as if randomized'' conditional on observed history:
\begin{equation}
\bigl(Y^g_{t+1},\ldots,Y^g_{K+1}\bigr)
\ \perp\!\!\!\perp\ 
\bigl(A_t, C_{t+1}\bigr)
\ \Bigm|\ 
\bar L_t=\bar l_t,\ \bar A_{t-1}=\bar a^g_{t-1},\ \bar C_t=\bar Y_t=0.
\end{equation}
This conditional independence concerns the assignment mechanisms for $A_t$ and
$C_{t+1}$, not the causal effect of treatment. It rules out unmeasured factors that,
given the observed history, jointly influence treatment decisions or censoring and the
future counterfactual outcomes under $g$.

\textbf{(2) Positivity (overlap).}
For histories that occur with positive density under the observed data-generating
process, there must be empirical support for the treatment levels prescribed by $g$
\emph{and} remaining uncensored. Because $A_t$ is continuous, we express positivity in
terms of neighborhoods. For each $k$ and each $\varepsilon>0$,
\begin{multline}
f_{\bar L_k,\bar A_{k-1},\bar C_k,\bar Y_k}\bigl(\bar l_k,\bar a_{k-1},0,0\bigr)>0
\ \Rightarrow\\
P\!\left(A_k \in \mathcal N_\varepsilon(a_k^g),\ C_{k+1}=0
\ \Bigm|\ \bar L_k=\bar l_k,\ \bar A_{k-1}=\bar a_{k-1},\ \bar C_k=\bar Y_k=0\right)>0,
\end{multline}
where $\mathcal N_\varepsilon(a)=\{a':|a'-a|<\varepsilon\}$ and
$a_k^g=g_k(\bar l_k,\bar a_{k-1})$. Intuitively, among histories that occur, the observed
data must contain nonzero probability of treatment values near those dictated by $g$,
while remaining alive and uncensored.

\textbf{(3) Consistency (well-defined interventions).}
If $\bar A_K=\bar a_K^g$ and $\bar C_{K+1}=0$, then $\bar L_K^g=\bar L_K$ and
$Y_{K+1}^g=Y_{K+1}$.

\subsubsection*{Plausibility of the positivity assumption}\label{sec:positivity}

Our target interventions are deterministic dynamic rules (closed-loop controllers) that
map histories to unique actions. In the extreme case where the observed treatment
mechanism followed a deterministic rule exactly, positivity would fail because for a
given history only one action would ever occur. In practice, ICU treatment exhibits
substantial stochastic variation across clinicians and time, including untreated
periods, intermittent dosing, and near--closed-loop titration. This heterogeneity
supports approximate overlap: within the range of histories observed, there is nonzero
probability of treatment values near those prescribed by reasonable target policies.

To further limit extrapolation, we restrict simulated interventions to operate within
empirically observed ranges of disease severity and dosing, thereby focusing estimation
on regions of the covariate--treatment space supported by the observed data.

Data generated according to the causal graph in Figure~\ref{fig:DAG2} satisfy these
identifiability conditions, as does the data-generating model considered in CICADAS.

\subsection{Identification via G-Formula for Causal Survival Analysis}

Under the above assumptions, the causal survival probability under policy $g$ is
identified by the parametric g-formula~\cite{robins1986new}. For our discrete-time survival outcome, the
probability of surviving through time $K+1$ under policy $g$ can be written as
\begin{align}
\label{eq:gformula_survival1}
P\bigl(Y^g_{K+1} = 0\bigr)
&= \sum_{\bar{l}_K}
\left[
\prod_{t=0}^{K}
P\Bigl(
Y_{t+1} = 0
\,\Bigm|\,
\bar{L}_t = \bar{l}_t,\,
\bar{A}_t = \bar{a}^g_t,\,
\bar{C}_t = \bar{Y}_t = 0
\Bigr)
\right]
f^g(\bar{l}_K),
\end{align}
where $f^g(\bar l_K)$ denotes the joint distribution of the covariate history $\bar L_K$
that would arise under policy $g$ under the hypothetical elimination of censoring (i.e.,
$\bar C=0$). Survival through time $K+1$ corresponds to the product over $t=0,\ldots,K$,
while the covariate history $\bar L_K$ is generated through transitions at
$t=0,\ldots,K-1$.

Although the estimand in \eqref{eq:gformula_survival1} is defined under elimination of
censoring, the observed data are subject to censoring. Assumption (1) is stated jointly
for treatment and censoring and therefore requires conditional independent censoring
given observed history; accordingly, we estimate the disease-evolution and outcome
models using only person-time during which individuals are alive and uncensored, without
fitting an explicit censoring model.

\paragraph{Simulating the covariate distribution under $g$.}
We do not estimate $f^g(\bar l_K)$ directly. Instead, we generate sample paths from it
by recursively simulating $L_{t+1}$ from the disease-evolution model
\[
f\bigl(L_{t+1}\mid \bar L_t,\bar A_t,\bar C_t=\bar Y_t=0\bigr)
\quad \text{with} \quad A_t = g_t(\bar L_t,\bar A_{t-1}),
\]
which yields trajectories distributed as $f^g(\bar L_K)$.
Under sequential ignorability and consistency, the conditional transition densities
under $g$ coincide with those estimable from observed data, giving the factorization
\[
f^g(\bar l_K)
=
\prod_{t=0}^{K-1}
f\!\left(
l_{t+1}
\,\middle|\,
\bar l_t,\,
\bar a_t^g,\,
\bar C_t=\bar Y_t=0
\right).
\]
A more detailed derivation is provided in Supplementary Section~S1.

\paragraph{Monte Carlo implementation.}
The summation in \eqref{eq:gformula_survival1} is generally intractable because it
requires integrating over all possible covariate trajectories. We therefore approximate
it by Monte Carlo simulation: generate $M$ sample paths
$\bar l_K^{(m)}=(l_0^{(m)},\ldots,l_K^{(m)})$ from the disease-evolution model under
policy $g$ (and $\bar C=0$), and compute
\begin{align}
P\bigl(Y^g_{K+1}=0\bigr)
&\approx
\frac{1}{M}\sum_{m=1}^{M}
\prod_{t=0}^{K}
P\Bigl(
Y_{t+1}=0
\,\Bigm|\,
\bar L_t=\bar l_t^{(m)},\,
\bar A_t=\bar a_t^g,\,
\bar C_t=\bar Y_t=0
\Bigr),
\end{align}
where $\bar l_t^{(m)}=(l_0^{(m)},\ldots,l_t^{(m)})$ is the covariate history along the
$m$-th simulated trajectory.

\paragraph{Models required for implementation.}
Implementing the Monte Carlo g-formula requires estimating two statistical models from
the observed (uncensored) data:
\begin{enumerate}
\item \textbf{Disease evolution model:}
\[
f\!\left(L_{t+1}\mid \bar L_t,\bar A_t,\bar C_t=\bar Y_t=0\right).
\]
\item \textbf{Outcome (mortality) model:}
\[
P\!\left(Y_{t+1}=1\mid \bar L_t,\bar A_t,\bar C_t=\bar Y_t=0\right).
\]
\end{enumerate}

\paragraph{Relation to alternative g-formula estimators.}
A comparison of the Monte Carlo parametric g-formula used here with iterated-conditional-expectation
and weighting-based alternatives is given in Supplementary Section~S1.3.
